\section{P-G assembly: uniform remainder, parameters, application of Parts 2--4, tori, measure, strictness}\label{bpP:p5}
\begingroup\sloppy
\providecommand{\wn}[3]{\lvert #1\rvert_{#2,#3}}

This part proves the frozen theorem \Pref{bpP:p0:pg} (Theorem~\ref{bpP:p5:main}) from Parts~1--4. Write $n=a+b$. The route is CP10 \S4, pp.~562--564, but the \emph{hypotheses are the frozen ones} and everything the source asserts without proof is supplied: the uniform remainder (CP10 (60)), the nondegeneracy of the frequency map on a complex neighbourhood, the explicit parameter choice with the numerical bounds of \Pref{bpP:p3:interface} tracked, the conversion of tori to \Lean{IsCylinderTorus}, the cell measure, and the strict inequalities.

\paragraph{What is different from CP10.} (1) Hypothesis (A$_2''$) is the Birkhoff form \emph{already}, so Steps~2--4 of CP10 \S3 (equilibrium, diagonalization, Birkhoff) are not needed: the averaged secular function $N$ of \Pref{bpP:p0:birkhoff} is used as it stands, and the remainder after averaging is $\mu Q+\widetilde P$, with $Q=P_{\rm av}-N=O(\abs{(p,q)}^5)$ and $\widetilde P$ of Part~2. (2) The KAM step is applied on the shell $\lambda_1<\xi_j<\lambda_2'$ with the constant $K=K_0$ \emph{bounded} (the two Diophantine scales $\gamma_1=\bar\gamma_2=\kappa_0 r^{5/2}$ make $K=\frac6s\log_+(\kappa_0^2/C^4)$ independent of $r$); this is why no power of $\log(1/r)$ enters the $\sqrt r$ loss of the measure. (3) The set of $\bar\omega_0$-Diophantine actions is a Cantor-type set $\check D$; its complex neighbourhood is not controlled by the collar, so Lemma~\ref{bpP:p5:budget} and Proposition~\ref{bpP:p5:lower} re-run the budget of \Pref{bpP:p4:budget}, \Pref{bpP:p4:cell-budget} for $D=\check D\times A_{\rm sh}$ (replacing $\bar D_t$ by $\check D\subset\bar D_t$ in the \emph{domain}, but not in the Diophantine count).

\subsection{Conventions and the order of the constants}\label{bpP:p5:conv}

All coordinates are CP's: angles $\varphi$ of period $2\pi$, actions $I$, oscillators $(p,q)$; the frozen fast frequency map is $\bar\omega_0=\omega_{H_0}=\partial H_0$ (CP angular convention). Part~3 works with period-one angles and $H_3=\tilde H(J,2\pi\theta)/2\pi$ (\Pref{bpP:p3:period-dictionary}); all quantities of Part~3 with $2\pi$ factors are absorbed in the constants $C$ below, and the dictionary $\gamma_i=\alpha_i/2\pi$ of \Pref{bpP:p4:freq} is used. Norms are max norms; $\wn{\cdot}{}{}$ is the weighted norm of \Nref{bpN:p4:norm} with coefficient radius (CP's $s$). $\lambda_0$ denotes the sup over real points of the complex bounds that appear.

\begin{remark}[Order in which every constant is chosen]\label{bpP:p5:order}
In the proof of Theorem~\ref{bpP:p5:main} the constants are chosen in this order; each depends only on the earlier ones.
\begin{enumerate}
\item \emph{Data and hypotheses} (fixed first): $a,b,V,W,A,d,r_0,\mu_0,B,H_0,P,P_0,\Omega,\beta$.
\item \emph{Analytic constants} (Lemmas~\ref{bpP:p5:ext}, \ref{bpP:p5:birkhoff-ext}, \ref{bpP:p5:freq-data}): $\delta_c$, then $s_*$ and $\varepsilon_{\rm an}=\min\{r_0/2,\delta_c/4\}$ (the complex oscillator radius is chosen after $\delta_c$), $r_1,\,\Lambda_P,\,C_Q,\,\rho_c,\,M_\beta,\,L_\beta,\,\Lambda_N,\,c_0,\,\bar M_0,\,R_y,\,L_g,\,\mu_*$ (data only; $\mu$ and $r$ do not enter).
\item $\tau:=n+1$; the Part~3 numerical constant $C_1\ge1$ of \Pref{bpP:p3:interface} (Lemma~\ref{bpP:p5:kam-hyp}); the widths $s,\bar s$ (functions of $s_*$).
\item The scale constants $\kappa_0$, then $\kappa_1$, then $G$ (Lemma~\ref{bpP:p5:kam-cond}, Definition~\ref{bpP:p5:params}); then the measure constant $C_*$ (Prop.~\ref{bpP:p5:lower}).
\item $r_*$ (finite list of smallness conditions, Definition~\ref{bpP:p5:rstar}); $C:=\max\{C_*+1,C_w\}$ with $C_w$ the frequency bound; $c:=2\tau+2=2n+4$.
\item Only \emph{then} $r\in(0,r_*)$ and $\mu\in(0,\min\{\mu_0,r^6/(\log\frac1r)^c\})$ are given; $K,\bar K,\bar\gamma,\bar\rho,\rho,\hat\rho,\hat E,r_K,t,\lambda_1,\lambda_2,\lambda_2'$ are explicit functions of $(r,\mu)$ (Definition~\ref{bpP:p5:params}); finally the set $S$ (with the deleted tori of \S\ref{bpP:p5:strict}).
\end{enumerate}
\end{remark}

\subsection{Uniform complex extension and uniform remainder}\label{bpP:p5:unif}

\emph{Warning (review P5-01).} Frozen real analyticity gives only a thin complex tube around the real box; it does \emph{not} give a holomorphic extension to the full complex oscillator polydisc of radius $r_0/2$ (a perturbation with a pole at $p=i\delta$, $0<\delta<r_0/4$, satisfies every frozen premise). The radius $\varepsilon_{\rm box}=r_0/2$ is used only for the \emph{real} box; the complex oscillator radius $\varepsilon_{\rm an}$ is chosen \emph{after} $\delta_c$ (Definition~\ref{bpP:p5:consts}), from the data only.

\begin{lemma}[Uniform holomorphic extension of the family]\label{bpP:p5:ext}
Assume (H1), (H2) of \Pref{bpP:p0:pg} and $\operatorname{HasFastDiffeo}(H_0,W)$. Put $\varepsilon_{\rm box}:=r_0/2$ (a \emph{real} box radius only) and, for $\delta>0$, let $\mathcal X_\delta\subset\C^a\times\C^a\times\C^b\times\C^b$ be the set of $(I,\varphi,p,q)$ with $\operatorname{dist}_\infty\bigl((I,\operatorname{Re}\varphi,\operatorname{Re}p,\operatorname{Re}q),\ \overline V\times\R^a\times\{\abs{(p,q)}_\infty\le\varepsilon_{\rm box}\}\bigr)<\delta$ and $\abs{\operatorname{Im}(I,\varphi,p,q)}_\infty<\delta$. Then there are $\delta_c\in(0,1)$ and $M_P<\infty$ and a map $\tilde P:[0,\mu_0]\times\mathcal X_{\delta_c}\to\C$ with: $\tilde P(\mu,\cdot)$ holomorphic, $\tilde P=P$ at real points, $\tilde P(\mu,\cdot+2\pi k e_{\rm f})=\tilde P(\mu,\cdot)$ ($e_{\rm f}$ the fast angles), $\abs{\tilde P}\le M_P$; $\tilde H_0$ holomorphic on $\{I:\operatorname{dist}_\infty(I,\overline V)<\delta_c\}$ with $\abs{\tilde H_0}$, $\norm{\partial\tilde H_0},\norm{\partial^2\tilde H_0}\le M_0$, and $\tilde P$ jointly continuous in $\mu$.
\end{lemma}
\begin{proof}
\proofstep{Local extensions}\label{stepP:p5:ext-loc} Let $X_0=[0,\mu_0]\times\overline V\times[0,2\pi]^a\times\{\abs{(p,q)}_\infty\le\varepsilon_{\rm box}\}$, compact. Since $2\varepsilon_{\rm box}^2<r_0^2$ for the sup-norm ball, $X_0\subset[0,\mu_0]\times\overline{\mathcal C(V,r_0)}$. By (H1) $P$ is analytic at each point of $X_0$: a convergent power series in a polydisc, which defines a holomorphic extension there. By compactness a Lebesgue number gives $\delta_1$ such that for every $x\in X_0$ the polydisc $B(x,\delta_1)$ carries a holomorphic extension $P_x$, with $\abs{P_x}\le M$ uniformly (shrink $\delta_1$ so that closed polydiscs lie in the domains of finitely many series).
\proofstep{Gluing}\label{stepP:p5:ext-glue} If $x,y\in X_0$ and $B(x,\delta_1)\cap B(y,\delta_1)\ne\emptyset$ then $\abs{x-y}<2\delta_1$ and the real midpoint lies in both; the intersection is convex hence connected, $P_x=P_y$ on its real points (both equal $P$), so $P_x=P_y$ on the intersection (identity theorem). Thus the $P_x$ define one holomorphic $\tilde P$ on $\bigcup_xB(x,\delta_1)$ (for each real $\mu$; jointly holomorphic in $(\mu,\cdot)$ near real $\mu$).
\proofstep{Periodicity}\label{stepP:p5:ext-per} By (H2) and continuity of $P$ on $\overline{\mathcal C(V,r_0)}$ (density of $\mathcal C(V,r_0)$ in its closure; $P$ analytic hence continuous) periodicity holds on the closure. For $\varphi$ outside $[0,2\pi]^a$ define $\tilde P$ by translation by $2\pi k_{\rm f}$ from the fundamental box; two such definitions agree on overlaps by the argument of Step~\ref{stepP:p5:ext-glue} (agreement at real points, convex intersection). This gives $\mathcal X_{\delta_c}$ for $\delta_c=\delta_1$.
\proofstep{$H_0$}\label{stepP:p5:ext-H0} $H_0$ is analytic on $W\supset\overline V$ (compact): Steps~\ref{stepP:p5:ext-loc}--\ref{stepP:p5:ext-glue} (without periodicity) give $\tilde H_0$ on a complex $\delta_c$-neighbourhood of $\overline V$ (decrease $\delta_c$); derivatives bounded by shrinking and Cauchy (\Kref{bp:p1:cauchy-open}).
\end{proof}

\begin{definition}[Strip, averaged function, constants]\label{bpP:p5:consts}
$s_*:=\min\{\delta_c,1\}/2$ (a coefficient radius $\le1/2$), $\varepsilon_{\rm an}:=\min\{\varepsilon_{\rm box},\delta_c/4\}$ (the \emph{complex oscillator radius}), $\hat D_c:=\{I\in\C^a:\operatorname{dist}_\infty(I,\overline V)<\delta_c/2\}$. For $(\mu,I,p,q)$ with $I\in\hat D_c$, $\abs{(p,q)}_\infty<\varepsilon_{\rm an}$ complex (such points, with $\abs{\operatorname{Im}\varphi}<s_*/6$, lie in $\mathcal X_{\delta_c}$: $\abs{\operatorname{Re}(p,q)}<\varepsilon_{\rm an}\le\varepsilon_{\rm box}$, $\abs{\operatorname{Im}(p,q)}<\delta_c/4$), put
\[\tilde P_{\rm av}(\mu,I;p,q)=(2\pi)^{-a}\int_{[0,2\pi]^a}\tilde P(\mu,I,\varphi,p,q)\,d\varphi\]
(real $\varphi$; $=P_{\rm av}$ of \Pref{bpP:p0:secavg} at real points of the cylinder). The N-norm of $\tilde P(\mu,\cdot)$ (Fourier in $\varphi$, coefficient radius $s_*$, sup over $(I,p,q)$ in the complex domain $\hat D_c\times B_{\varepsilon_{\rm an}}\times B_{\varepsilon_{\rm an}}$) satisfies
\[\Lambda_P:=\sup_{\mu\in[0,\mu_0]}\wn{\tilde P(\mu,\cdot)}{}{s_*}\le M_P\coth(\delta_c/4)^a<\infty\quad(\ge1\ \text{after enlarging}).\]
\end{definition}
\begin{proof}[Justification of finiteness]
$\abs{\hat f_k}\le\sup\abs f\,e^{-\sigma_1\abs k_1}$ on a strip of width $\sigma_1$ (Cauchy in the angle, $\hat f_k=(2\pi)^{-a}\int f(\varphi)e^{-ik\cdot\varphi}d\varphi$ shifted to $\operatorname{Im}\varphi=\mp\sigma_1\operatorname{sgn}k$), so $\sum_k\abs{\hat f_k}e^{s_*\abs k_1}\le\sup\abs f\sum_ke^{-(\delta_c-s_*)\abs k_1}\le M_P\coth(\delta_c/4)^a$ since $\delta_c-s_*\ge\delta_c/2$. This is the elementary \emph{sup-to-weighted conversion}, used again in Lemma~\ref{bpP:p5:kam-hyp} (cf.\ \Nref{bpN:p4:sup-weighted}, \Kref{bp:p1:decay}).
\end{proof}

\begin{lemma}[Birkhoff coefficients as derivatives; analytic extension to $\overline V$]\label{bpP:p5:birkhoff-ext}
Assume (H1), (H2), (H4), (H5), (H6) of \Pref{bpP:p0:pg} ((H5) is used only for (c)). For $(\mu,I)\in[0,\mu_0]\times\hat D_c$ define, with the holomorphic extension $\tilde P_{\rm av}$ (and $R_j=(p_j^2+q_j^2)/2$)
\[\tilde P_0=\tilde P_{\rm av}\big|_{x=0},\quad\tilde\Omega_j=\partial_{p_j}^2\tilde P_{\rm av}\big|_{x=0},\quad\tilde\beta_{ij}=\partial^2_{p_i}\partial^2_{p_j}\tilde P_{\rm av}\big|_{x=0}\ (i\ne j),\quad\tilde\beta_{jj}=\tfrac13\partial^4_{p_j}\tilde P_{\rm av}\big|_{x=0},\]
$x=(p,q)$, and $\tilde N=\tilde P_0+\tilde\Omega\cdot R+\frac12R\cdot\tilde\beta R$, $Q:=\tilde P_{\rm av}-\tilde N$. Then:
(a) $\tilde P_0,\tilde\Omega,\tilde\beta$ are holomorphic in $I\in\hat D_c$, continuous in $\mu\in[0,\mu_0]$, $\tilde\beta$ is symmetric, and they \emph{equal} $P_0,\Omega,\beta$ on $[0,\mu_0]\times V$;
(b) for each $\mu\in[0,\mu_0]$ and $I\in\hat D_c$, the jet of $Q(\mu,I;\cdot)$ at $x=0$ vanishes to order $4$ (all derivatives of order $\le4$);
(c) $\abs{\det\tilde\beta(\mu,I)}\ge B$ for all $(\mu,I)\in[0,\mu_0]\times\overline V$.
\end{lemma}
\begin{proof}
\proofstep{Coefficients on $V$}\label{stepP:p5:be-coef} Fix $\mu\in[0,\mu_0]$, $I\in V$. By \Pref{bpP:p0:birkhoff}(ii), $P_{\rm av}-N=o(\abs x^4)$, and $P_{\rm av}-N$ is real analytic near $x=0$ (analytic by Lemma~\ref{bpP:p5:ext} and differentiation under the integral; $N$ polynomial). Hence the Taylor polynomial of $P_{\rm av}(\mu,I;\cdot)$ of degree $\le4$ equals that of $N$, which is $N$. Compute from $R_j=(p_j^2+q_j^2)/2$: $\partial^2_{p_j}N(0)=\Omega_j$; $R_iR_j=\frac14(p_i^2+q_i^2)(p_j^2+q_j^2)$ gives $\partial^2_{p_i}\partial^2_{p_j}N(0)=\frac12(\beta_{ij}+\beta_{ji})=\beta_{ij}$ for $i\ne j$ \emph{using} (H6); $\frac12\beta_{jj}R_j^2=\frac18\beta_{jj}(p_j^2+q_j^2)^2$ gives $\partial^4_{p_j}N(0)=3\beta_{jj}$; $P_0=P_{\rm av}|_{x=0}$. So (a) on $V$.
\proofstep{Analyticity}\label{stepP:p5:be-an} The listed derivatives of $\tilde P_{\rm av}$ at $x=0$ are holomorphic in $I$ (integral of holomorphic functions, uniform bounds by Lemma~\ref{bpP:p5:ext}), continuous in $\mu$; symmetry of $\tilde\beta$ is symmetry of mixed derivatives.
\proofstep{(b) on $V$, then everywhere}\label{stepP:p5:be-jet} On $V$, $\tilde N=N$ and Step~\ref{stepP:p5:be-coef} give the vanishing jet. Each derivative $\partial_x^\alpha Q(\mu,\cdot;0)$, $\abs\alpha\le4$, is holomorphic on $\hat D_c$, which is connected ($\overline V$ is connected as the closure of the connected $V$ (\texttt{hVc}), and $\hat D_c$ is a union of balls centered on it) and vanishes on the real open set $V$; by the identity theorem it vanishes on $\hat D_c$.
\proofstep{(c)}\label{stepP:p5:be-det} On $V$ by (H5) and (a); $\overline V$ by continuity of $\det\tilde\beta$ in $I$ and $V$ dense in $\overline V$.
\end{proof}

\begin{proposition}[Uniform remainder: CP10 (60) from the frozen hypotheses]\label{bpP:p5:uniform-remainder}
Under the hypotheses of Lemma~\ref{bpP:p5:birkhoff-ext} (that is, (H1), (H2), (H4)--(H6)) there are $r_1\in(0,\varepsilon_{\rm an}]$ and $C_Q<\infty$ (depending on the data only) such that for all $\mu\in[0,\mu_0]$, $I\in\hat D_c'$ ($:=\{\operatorname{dist}_\infty(I,\overline V)<\delta_c/4\}$) and complex $x=(p,q)$ with $\abs x_\infty\le r_1$:
\[\abs{Q(\mu,I;x)}\le C_Q\abs x_\infty^5.\]
In particular $\abs{P_{\rm av}-\tilde N}\le C_Q\abs{(p,q)}_\infty^5$ on $[0,\mu_0]\times\overline V\times\{\abs x_\infty\le r_1\}$ (real), where $\tilde N,\tilde P_0,\tilde\Omega,\tilde\beta$ are the \emph{constructed extensions} of Lemma~\ref{bpP:p5:birkhoff-ext}, analytic on $\hat D_c$, with $\abs{\det\tilde\beta}\ge B$ on $\overline V$. They equal the original $P_0,\Omega,\beta,N$ \emph{on $V$ only} (the frozen total functions are unconstrained on $\partial V$ and outside, so no statement about their boundary values is made); on $V$ the bound reads $\abs{P_{\rm av}-N}\le C_Q\abs{(p,q)}^5_\infty$. \emph{Hypothesis (A$_2''$) with a $\delta$ depending on $(\mu,I,\epsilon)$ is thus upgraded to a uniform $O(\abs x^5)$ bound.}
\source{CP10 (60), p.~563, asserted ``for a suitable $C>0$'' from (A$_2''$); proof not given.}
\sourcegap{CP10 (60) is asserted without the compactness and analyticity argument; supplied here.}
\end{proposition}
\begin{proof}
\proofstep{Bound}\label{stepP:p5:ur-bound} Take $r_1:=\varepsilon_{\rm an}/2$. $Q$ is holomorphic in $x$ on $\{\abs x_\infty<\varepsilon_{\rm an}\}$ and $\abs Q\le M_P+\abs{\tilde N}\le M_1$ on $\hat D_c'\times\{\abs x_\infty\le r_1\}$, $r_1=\varepsilon_{\rm an}/2$, with $M_1$ from the bounds of Lemma~\ref{bpP:p5:ext} and Cauchy estimates for $\tilde\Omega,\tilde\beta$ (derivatives at $0$ of a function bounded by $M_P$ on a polydisc of radius $\varepsilon_{\rm an}$).
\proofstep{Schwarz}\label{stepP:p5:ur-schwarz} Fix $(\mu,I)$ and $\hat x$ with $\abs{\hat x}_\infty=1$; $g(t)=Q(\mu,I;t\hat x)$ is holomorphic on $\abs t<r_1$, $\abs g\le M_1$, and by Lemma~\ref{bpP:p5:birkhoff-ext}(b) $g$ vanishes to order $\ge5$ at $0$ (derivatives of $g$ at $0$ of order $\le4$ are combinations of derivatives of $Q$ of order $\le4$). So $g(t)/t^5$ is holomorphic and bounded by $M_1/r_1^5$ on the boundary circles, hence on the disc (maximum principle): $\abs{g(t)}\le M_1\abs t^5/r_1^5$. Take $C_Q=M_1r_1^{-5}$.
\end{proof}

\subsection{The fast frequency map on a complex neighbourhood; the nondegenerate frequency map of the averaged system (domain reduction, CP10 (58))}\label{bpP:p5:domain-reduction}

The frozen hypothesis $\operatorname{HasFastDiffeo}(H_0,W)$ replaces CP10's assumption that $\bar\omega_0=\partial H_0$ is a diffeomorphism of $V_{\rho_0}$ with bounded inverse Jacobian (58). Write $Z_c:=\hat D_c$-type neighbourhoods $X_\rho:=\{z\in\C^N:\operatorname{dist}_\infty(z,X)<\rho\}$ for $X\subset\R^N$.

\begin{lemma}[Lipschitz bound on a complex neighbourhood]\label{bpP:p5:lip-complex}
Let $X\subset\R^N$, $\rho>0$, $\Delta:X_{2\rho}\to\C^m$ holomorphic with $\sup\abs\Delta\le A$. Then $\abs{\Delta(z)-\Delta(z')}\le\frac{2NA}{\rho}\abs{z-z'}$ for all $z,z'\in X_\rho$.
\end{lemma}
\begin{proof}
Close pairs $\abs{z-z'}<\rho$: every point $w$ of the segment $[z,z']$ is within $\abs{z-z'}/2<\rho/2$ of the nearer endpoint, which lies in $X_\rho$, so $w\in X_{3\rho/2}$ and the polydisc $B(w,\rho/2)\subset X_{2\rho}$; Cauchy (\Kref{bp:p1:cauchy-open}) gives $\norm{D\Delta(w)}\le2NA/\rho$; mean value inequality (\Kref{bp:p1:ext-mvt}). Distant pairs: $\abs{\Delta(z)-\Delta(z')}\le2A\le(2A/\rho)\abs{z-z'}$.
\source{Ours; the complex-point version of \Pref{bpP:p3:lip-sup}.}
\end{proof}

\begin{lemma}[Complex bi-Lipschitz bound for $\bar\omega_0$]\label{bpP:p5:bilip}
Under $\operatorname{HasFastDiffeo}(H_0,W)$ there are $\rho_c\in(0,\delta_c/16)$ and $c_0>0$ with $\abs{\tilde\omega_0(z)-\tilde\omega_0(z')}\ge c_0\abs{z-z'}$ for all $z,z'\in(\overline V)_{\rho_c}$, where $\tilde\omega_0=\partial\tilde H_0$; hence $\norm{(\partial\tilde\omega_0(z))^{-1}}\le c_0^{-1}$ there. Moreover $\bar M_0:=\sup_{(\overline V)_{2\rho_c}}\norm{\partial\tilde\omega_0}<\infty$.
\end{lemma}
\begin{proof}
\proofstep{Real facts}\label{stepP:p5:bl-real} By \Pref{bpP:p0:fastdiffeo}, $g(\bar\omega_0(I))=I$ on $W$, so $\bar\omega_0$ is injective on $W\supset\overline V$; differentiating (chain rule, $g\in C^1$, $\bar\omega_0$ analytic) $Dg\,D\bar\omega_0=1$, so $\partial^2H_0=D\bar\omega_0$ is invertible on $W$.
\proofstep{Compactness}\label{stepP:p5:bl-comp} Suppose not: $z_m\ne z_m'\in(\overline V)_{1/m}$ with $\abs{\tilde\omega_0(z_m)-\tilde\omega_0(z_m')}<\abs{z_m-z_m'}/m$. Pass to subsequences with $z_m\to x$, $z_m'\to x'\in\overline V$ (compact). If $x\ne x'$, continuity gives $\bar\omega_0(x)=\bar\omega_0(x')$, contradicting injectivity. If $x=x'$: let $m_0=\norm{(\partial\tilde\omega_0(x))^{-1}}$; choose $\delta$ with $\norm{\partial\tilde\omega_0(w)-\partial\tilde\omega_0(x)}\le1/(2m_0)$ on $B(x,\delta)$ (continuity); for $z,z'\in B(x,\delta)$, $\abs{\tilde\omega_0(z)-\tilde\omega_0(z')}\ge\abs{\partial\tilde\omega_0(x)(z-z')}-\frac1{2m_0}\abs{z-z'}\ge\frac1{2m_0}\abs{z-z'}$, contradicting the choice for $m>2m_0$. The operator bound: differentiate along $z'\to z$. $\bar M_0<\infty$ by Lemma~\ref{bpP:p5:ext}.
\end{proof}

\begin{lemma}[One-scale Diophantine measure on $V$ (CP10 (38))]\label{bpP:p5:dio-one}
Let $L_g,R_y,\Lambda_g$ be as in \Pref{bpP:p4:lip-inverse} and $S_1$ as in \Pref{bpP:p4:lattice-tail} (with $n=a+b$, $\tau>n$). For $\bar\gamma>0$ put $\mathcal D^1_{\bar\gamma}:=\{y\in\R^a:\abs{k\cdot y}\ge\bar\gamma\abs k_1^{-\tau}\ \forall k\in\Z^a\setminus0\}$ (closed). Then $\meas\bigl(V\setminus\bar\omega_0^{-1}(\mathcal D^1_{\bar\gamma})\bigr)\le c_{\rm dio}\bar\gamma$ with $c_{\rm dio}:=2\Lambda_g^aS_1(2R_y)^{a-1}$.
\end{lemma}
\begin{proof}
$V\setminus\bar\omega_0^{-1}(\mathcal D^1_{\bar\gamma})\subset\bigcup_kg(Y_k)$, $Y_k=\{y\in Y_0:\abs{k\cdot y}<\bar\gamma\abs k_1^{-\tau}\}$ (a point $I\in V$ with $\bar\omega_0(I)\notin\mathcal D^1_{\bar\gamma}$ has $\bar\omega_0(I)\in Y_k$ for some $k$ and $I=g(\bar\omega_0(I))$, \Pref{bpP:p4:lip-inverse}; containment suffices). By \Pref{bpP:p4:affine-slab} (with $c=0$, $N=a$, $\lambda=R_y$, $\ell=k$, $\abs\ell\ge1$), $\meas Y_k\le2\bar\gamma\abs k_1^{-\tau}(2R_y)^{a-1}$; $g$ is $\Lambda_g$-Lipschitz on $Y_0$, so \Kref{bp:p1:lip-image} (after McShane extension \Kref{bp:p1:mcshane}) gives $\meas g(Y_k)\le\Lambda_g^a\meas Y_k$. Sum over $k$ using $\sum_{k\in\Z^a\setminus0}\abs k_1^{-\tau}\le S_1$ (the lattice $\Z^a\subset\Z^n$).
\source{CP10 (38), p.~556, ``standard measure estimate''; proof ours.}
\end{proof}

\begin{lemma}[Constants of the averaged frequency map]\label{bpP:p5:freq-data}
Let $\tilde N$ be as in Lemma~\ref{bpP:p5:birkhoff-ext}. There are constants $A_N,\Lambda_N,M_\beta,L_\beta\ge1$ (data only) such that on $\mathcal E:=(\overline V)_{\rho_c}\times\{\abs{J_2}_\infty<2\}$, $\tilde N$, $\partial_{J}\tilde N$, $\partial^2_J\tilde N$ are bounded by $A_N$; with $G:=\partial_{J_1}\tilde N$ and $\Theta(J_1,\xi):=\tilde\Omega(J_1)+\tilde\beta(J_1)\xi$, for $J_1,J_1'\in(\overline V)_{\rho_c}$, $\abs{\xi},\abs{\xi'}\le1$:
\[\abs{G(J_1,\xi)-G(J_1',\xi)},\ \abs{\Theta(J_1,\xi)-\Theta(J_1',\xi)}\le\Lambda_N\abs{J_1-J_1'},\quad\abs{G(J_1,\xi)-G(J_1,\xi')}\le\Lambda_N\abs{\xi-\xi'},\]
and $\abs\nu\le L_\beta\abs{\tilde\beta(J_1)\nu}$ for $J_1\in(\overline V)_{\rho_c}$, $\nu\in\C^b$ (after decreasing $\rho_c$). The real-variable constants of \Pref{bpP:p4:data} ((N1)--(N3) with $\bar D=\bar D_t$, $\lambda_0=1$) hold with the same $\Lambda_N,L_\beta$ and for all $\mu\in(0,\mu_0]$. (On $\bar D_t\subset V$ the extensions $\tilde P_0,\tilde\Omega,\tilde\beta$ coincide with the frozen $P_0,\Omega,\beta$, Lemma~\ref{bpP:p5:birkhoff-ext}(a), so Part~4 is applied to the original functions only at points of $V$.)
\end{lemma}
\begin{proof}
\proofstep{Bounds}\label{stepP:p5:fd-bd} $\tilde N$ is a polynomial in $J_2$ with coefficients $\tilde P_0,\tilde\Omega,\tilde\beta$ holomorphic and bounded on $\hat D_c$ (Lemma~\ref{bpP:p5:birkhoff-ext}, Cauchy from $M_P$); derivatives in $J_1$ up to order $3$ are bounded on $(\overline V)_{2\rho_c}$ for $\rho_c\le\delta_c/16$.
\proofstep{Lipschitz}\label{stepP:p5:fd-lip} In $J_1$: Lemma~\ref{bpP:p5:lip-complex} applied to $\partial_{J_1}\tilde N(\cdot,\xi)$ and $\Theta(\cdot,\xi)$ ($\abs\xi\le1$ makes the sup bound uniform). In $J_2$: polynomial of degree $\le2$ with bounded coefficients on the convex set $\abs\xi\le1$.
\proofstep{Inverse of $\tilde\beta$}\label{stepP:p5:fd-beta} $M_\beta:=\sup\abs{\tilde\beta_{ij}}$ on $(\overline V)_{\rho_c}$; $\abs{\det\tilde\beta}\ge B$ on $\overline V$ (Lemma~\ref{bpP:p5:birkhoff-ext}(c)), and $\det$ is Lipschitz, so $\abs{\det\tilde\beta}\ge B/2$ on $(\overline V)_{\rho_c}$ after decreasing $\rho_c$. Cramer as in \Pref{bpP:p4:beta-inverse} (the adjugate bound holds for complex matrices): $L_\beta:=2b!M_\beta^{b-1}/B$ (and $\tilde\beta^{\top}=\tilde\beta$).
\end{proof}

\begin{lemma}[Nondegeneracy and injectivity of the frequency map on the complex neighbourhood]\label{bpP:p5:freq-inj}
For $\mu\in(0,\mu_*]$, where $\mu_*\le1$ is defined by
\[\mu_*A_N\le\tfrac{c_0}2,\qquad 4\mu_*A_N^2L_\beta\le c_0,\qquad2\mu_*\Lambda_N^2L_\beta<c_0,\]
let $\tilde h_\mu(J)=\tilde H_0(J_1)+\mu\tilde N(J)$ on $\mathcal E_1:=(\overline V)_{\rho_c}\times\{\abs{J_2}<1\}$ and $F_\mu=\partial\tilde h_\mu=\bigl(\tilde\omega_0(J_1)+\mu G(J),\ \mu\Theta(J_1,J_2)\bigr)$, $U=\partial^2\tilde h_\mu$. Then
\begin{enumerate}[label=(\alph*)]
\item $F_\mu$ is injective on $\mathcal E_1$;
\item $U$ is invertible on $\mathcal E_1$, and with $U=\begin{pmatrix}U_{11}&U_{12}\\U_{21}&U_{22}\end{pmatrix}$ ($a+b$ blocks), $\hat U=(U_{21},U_{22})$, $U^{-1}=\binom{T_1}{T_2}$:
\[\norm U\le C_U,\quad\norm{\hat U}\le C_U\mu,\quad\norm{U^{-1}}\le C_U\mu^{-1},\quad\norm{T_1}\le C_U,\quad\norm{T_2}\le C_U\mu^{-1},\]
with $C_U:=\max\{\bar M_0+A_N+1,\ 4(A_N+1)^2L_\beta c_0^{-1}+4,\ 8L_\beta+4\}$ (depends on data only).
\end{enumerate}
\end{lemma}
\begin{proof}
\proofstep{(a)}\label{stepP:p5:fi-inj} Let $F_\mu(J)=F_\mu(J')$. First components: $c_0\abs{J_1-J_1'}\le\abs{\tilde\omega_0(J_1)-\tilde\omega_0(J_1')}=\mu\abs{G(J)-G(J')}\le\mu\Lambda_N(\abs{J_1-J_1'}+\abs{J_2-J_2'})$ (Lemma~\ref{bpP:p5:bilip}, Lemma~\ref{bpP:p5:freq-data}, splitting $G(J_1,J_2)-G(J_1',J_2')$ through $G(J_1',J_2)$). Since $\mu\Lambda_N\le c_0/2$ ($\Lambda_N\le A_N$ w.l.o.g.), $\abs{J_1-J_1'}\le(2\mu\Lambda_N/c_0)\abs{J_2-J_2'}$. Second components: $\tilde\beta(J_1)(J_2-J_2')=-[\Theta(J_1,J_2')-\Theta(J_1',J_2')]$, so $\abs{J_2-J_2'}\le L_\beta\Lambda_N\abs{J_1-J_1'}\le(2\mu\Lambda_N^2L_\beta/c_0)\abs{J_2-J_2'}$, hence $J_2=J_2'$, then $J_1=J_1'$.
\proofstep{(b) Blocks}\label{stepP:p5:fi-blk} $U_{11}=\partial\tilde\omega_0+\mu\tilde N_{11}=:A$, $U_{12}=\mu\tilde N_{12}$, $U_{21}=\mu\tilde N_{21}=U_{12}^\top$, $U_{22}=\mu\tilde\beta$. $\norm{A^{-1}}\le2/c_0$ (Lemma~\ref{bpP:p5:bilip}; Neumann series since $\mu A_N\le c_0/2$). Schur complement $S=\mu(\tilde\beta-\mu\tilde N_{21}A^{-1}\tilde N_{12})$; as $\mu A_N^2\cdot\frac2{c_0}L_\beta\le\frac12$, $\norm{\tilde\beta^{-1}}\le L_\beta$ gives $\norm{S^{-1}}\le2L_\beta/\mu$.
\proofstep{Inverse}\label{stepP:p5:fi-inv} Block inversion: $T_{11}=A^{-1}+\mu^2A^{-1}\tilde N_{12}S^{-1}\tilde N_{21}A^{-1}$, $T_{12}=-\mu A^{-1}\tilde N_{12}S^{-1}$, $T_{21}=-\mu S^{-1}\tilde N_{21}A^{-1}$, $T_{22}=S^{-1}$. Their norms are $\le2/c_0+8\mu A_N^2L_\beta c_0^{-2}$, $4A_NL_\beta/c_0$, $4A_NL_\beta/c_0$, $2L_\beta/\mu$. Row sums of $T_1=(T_{11},T_{12})$ are $O(1)$, of $T_2$ are $O(\mu^{-1})$, and the entries of $\hat U$ are $\mu\times$ bounded. The constant $C_U$ dominates all of these.
\end{proof}
\begin{remark}\label{bpP:p5:freq-inj-rem}
(a) gives injectivity on the \emph{complex} neighbourhood that \Pref{bpP:p3:interface} lists as an obligation of Part~5 (``injectivity/nondegeneracy of $\omega_\mu$ on a complex neighbourhood''); it uses only $\operatorname{HasFastDiffeo}$ (through Lemma~\ref{bpP:p5:bilip}), joint analyticity and (H5). The same computation, with $H_3=\tilde H(\cdot,2\pi\cdot)/2\pi$, gives the bounds for Part~3 with every number divided/multiplied by $2\pi$ as in \Pref{bpP:p3:period-dictionary}; those factors are absorbed in $C_1$ of Lemma~\ref{bpP:p5:kam-hyp}.
\end{remark}

\subsection{The parameters}\label{bpP:p5:parameters}

\begin{definition}[Parameters as functions of $(r,\mu)$]\label{bpP:p5:params}
Let $\tau:=n+1$ and $c:=2\tau+2$. Fix the data constants of Remark~\ref{bpP:p5:order}\,(2), $\sigma_{\rm av}:=s_*/6$, $\sigma_{\rm tot}:=s_*/12$, $s:=\sigma_{\rm tot}/5$, $\bar s:=4\sigma_{\rm tot}/5$ (so $4s=\bar s<1$ and $\bar s+s=\sigma_{\rm tot}$), and $C_1$, $\kappa_0$, $\kappa_1$, $G$ as in Lemmas~\ref{bpP:p5:kam-hyp}, \ref{bpP:p5:kam-cond}. For $r\in(0,1)$ and $\mu\in(0,\mu_0)$ put ($\mu_2:=\Lambda_P\mu$)
\begin{align*}
&\bar K:=\tfrac{30}{s_*}\log\tfrac1r,\quad t:=\sqrt r,\quad\lambda_2:=\tfrac{r^2}2,\quad\lambda_1:=\kappa_1r^{5/2},\quad\lambda_2':=\lambda_2-2\lambda_1,\quad\ell:=\lambda_2'-\lambda_1=\lambda_2-3\lambda_1,\quad\rho:=\tfrac{\lambda_1}2,\\
&\bar\gamma:=G\max\Bigl\{\sqrt{\mu_2}\,\bar K^{\tau+\frac12}r^{-5/2},\ r^{5/2}\bar K^{\tau+1}\Bigr\},\quad\alpha_{\rm av}:=\frac{\bar\gamma}{2\bar K^\tau},\quad\bar\rho:=\frac{\bar\gamma}{2\bar M_0\bar K^{\tau+1}},\\
&\gamma_1:=\kappa_0r^{5/2},\quad\gamma_2:=\mu\kappa_0r^{5/2},\quad K:=K_0:=\tfrac6s\log_+(\kappa_0^2/C_1^4),
\end{align*}
$\hat\rho,\hat E,r_K$ as in \Pref{bpP:p3:constants} with the numerical bounds of Lemma~\ref{bpP:p5:kam-hyp}. The sets: $\bar D_t$ of \Pref{bpP:p4:s-collar}, $\check D:=\bar D_t\cap\bar\omega_0^{-1}(\mathcal D^1_{\bar\gamma})$ (compact: $\mathcal D^1_{\bar\gamma}$ is closed, $\bar\omega_0$ continuous on $W\supset\bar D_t$), $A_{\rm sh}:=(\lambda_1,\lambda_2')^b$, $D:=\check D\times A_{\rm sh}$. The $2\pi$-dictionary: $\alpha_i:=2\pi\gamma_i$.
\end{definition}

\begin{definition}[Smallness thresholds]\label{bpP:p5:rstar}
$r_*\in(0,\min\{r_0,e^{-1/5},1\})$ is so small that for $0<r<r_*$ (all conditions involve only the constants of Remark~\ref{bpP:p5:order}, items 2--4, and $C_*$ of Proposition~\ref{bpP:p5:lower}):
\begin{enumerate}[label=(r\arabic*)]
\item $\log\frac1{r_*}\ge\Lambda_P$ and $r_*^2(\log\frac1{r_*})^{\tau+1}\le1$, so that $\mu<r^6/(\log\frac1r)^c$ implies $\mu_2\le r^6(\log\frac1r)^{-(2\tau+1)}$; and $r_*^6\le\mu_1:=\min\{\mu_*,\ (2(1+M)a\Lambda_N\Lambda_g)^{-1}\}$, $M=2b\Lambda_NL_\beta$ (\Pref{bpP:p4:data});
\item $\sqrt{r_*}<d$, $A\sqrt{r_*}\le\frac12\meas V$ (so $\meas\bar D_t\ge\frac12\meas V$: positive inner core), $r_*\le r_1/2$, $r_*\le\varepsilon_{\rm an}/4$;
\item with $C_\gamma:=G(30/s_*)^{\tau+1}$: $C_\gamma\sqrt{r_*}\le2\bar M_0\rho_c$ and $c_{\rm dio}C_\gamma\sqrt{r_*}\le\frac14\meas V$;
\item $6\kappa_1\sqrt{r_*}\le\frac12$;
\item the displacement margin $\sqrt2\,\Delta_{\max}\le\kappa_1r^{3/2}$ of Lemma~\ref{bpP:p5:location}, and $(C_*+1)\sqrt{r_*}\le\frac12$.
\end{enumerate}
\end{definition}

\begin{lemma}[Consequences of the parameter choice]\label{bpP:p5:param-facts}
For $0<r<r_*$, $0<\mu<\min\{\mu_0,r^6/(\log\frac1r)^c\}$:
(a) $\mu_2\le r^6(\log\frac1r)^{-(2\tau+1)}$; $\mu<\mu_1$; $\bar K s_*=30\log\frac1r\ge6$; $\bar\gamma\le C_\gamma\sqrt r$.
(b) $\bar\rho\le\rho_c$, $\bar\rho\ge\lambda_1$ if $G\ge2\bar M_0\kappa_1$ (so $\rho=\lambda_1/2\le\bar\rho/2$), and $\ell\ge\lambda_2/2\ge r^2/4$.
(c) $\meas\bar D_t\ge\frac12\meas V$, and $\meas\check D\ge\frac14\meas V$, $\meas\bar D_t\le4\meas\check D$.
\end{lemma}
\begin{proof}
(a) $\mu<r^6(\log\frac1r)^{-c}$ and $\Lambda_P\le\log\frac1r$ give the first; $\mu_2\bar K^{2\tau+1}\le r^6(30/s_*)^{2\tau+1}$, so $\sqrt{\mu_2}\bar K^{\tau+1/2}r^{-5/2}\le(30/s_*)^{\tau+1/2}\sqrt r$, and $r^{5/2}\bar K^{\tau+1}\le(30/s_*)^{\tau+1}r^{1/2}(r^2\log^{\tau+1}\frac1r)$, both $\le(30/s_*)^{\tau+1}\sqrt r$ by (r1). (b) $\bar\rho\le\bar\gamma/(2\bar M_0)\le C_\gamma\sqrt r/(2\bar M_0)\le\rho_c$ by (r3); $\bar\rho\ge G r^{5/2}/(2\bar M_0)\ge\lambda_1$ from the second branch of $\bar\gamma$; $\ell=\lambda_2(1-6\kappa_1\sqrt r)\ge\lambda_2/2$ by (r4). (c) \Pref{bpP:p4:collar}(b): $\meas\bar D_t\ge\meas V-At\ge\frac12\meas V$ (r2); $\meas(\bar D_t\setminus\check D)\le\meas(V\setminus\bar\omega_0^{-1}\mathcal D^1_{\bar\gamma})\le c_{\rm dio}\bar\gamma\le\frac14\meas V$ (Lemma~\ref{bpP:p5:dio-one}, (r3)); here $\frac12-\frac14=\frac14$.
\end{proof}

\subsection{Step 1 of CP10 \S4: averaging (Part 2)}\label{bpP:p5:kam-application}

\begin{proposition}[Averaged system in CP variables]\label{bpP:p5:avg}
For $r,\mu$ as above there is a map $\bar\Psi$ on $\mathcal W:=\{(I,\varphi,p,q)\in\C^{2n}:\ I\in\check D_{\bar\rho/2},\ \abs{\operatorname{Im}\varphi}<\sigma_{\rm av},\ \abs{(p,q)}_\infty<\varepsilon_{\rm an}/2\}$ (angles of period $2\pi$), holomorphic, injective, real at real points, with $\bar\Psi-\mathrm{id}$ $2\pi$-periodic in $\varphi$, mapping $\mathcal W$ into the domain of $\tilde P$ of Lemma~\ref{bpP:p5:ext}, with $D\bar\Psi^\top M_0D\bar\Psi=M_0$ (it preserves the frozen $\Omega$; this is the symplectic conclusion of \Pref{bpP:p2:osc-averaging} in the frozen order; the condition is unchanged if the form is multiplied by $-1$), and
\[H_\mu\circ\bar\Psi=\tilde H_0(I)+\mu\tilde N(I,R)+\mu Q(\mu,I;p,q)+\widetilde P_{\rm rem},\qquad\wn{\widetilde P_{\rm rem}}{}{\sigma_{\rm av}}\le2\Lambda_P\mu r^5\quad(\text{on }\mathcal W),\]
and, writing $(\bar I,\bar\varphi,\bar p,\bar q)\mapsto(I,\varphi,p,q)$: $\abs{I-\bar I}\le\Delta_I:=\dfrac{12\mu_2\bar K^\tau}{s_*\bar\gamma}$, $\abs{p-\bar p},\abs{q-\bar q}\le\Delta_{\max}:=\dfrac{18\mu_2\bar K^\tau}{\varepsilon_{\rm an}\bar\gamma}$, where, by the choice of $\bar\gamma$, $\Delta_{\max}\le\dfrac{18}{G\varepsilon_{\rm an}}r^{11/2}$, $\Delta_I\le\dfrac{12}{Gs_*}r^{11/2}$.
\end{proposition}
\begin{proof}
\proofstep{Dictionary}\label{stepP:p5:av-dict} Apply \Pref{bpP:p2:averaged-remainder} with $n=a$, $m=b$, $\bar D=\check D$, $\bar\rho=\bar\rho_{\rm CP}$, $\epsilon=r$, $\mu\to\mu_2$, $P\to\tilde P/\Lambda_P$ (so $\wn{\cdot}{}{}\le1$ by Definition~\ref{bpP:p5:consts}, and $H_\mu=H_0+\mu_2(\tilde P/\Lambda_P)$), in CP variables through the canonical chart $(\theta,I,p,q)=(\varphi/2\pi,2\pi I_{\rm CP},p,q)$ of \Pref{bpP:p2:use5}\,(3): $s_0=s_*/2\pi$, $\bar\alpha=2\pi\alpha=\alpha_{\rm av}$, $d=\min\{\bar\rho s_*,\varepsilon_{\rm an}^2\}=\bar\rho s_*$ ($\bar\rho\le\rho_c<\varepsilon_{\rm an}^2/s_*$ after decreasing $\rho_c$), $\kappa=\bar Ks_*=30\log\frac1r$ and $e^{-\kappa/6}=r^5$.
\proofstep{(R1)}\label{stepP:p5:av-R1} Let $I'\in\check D_{\bar\rho}$, $I\in\check D$ with $\abs{I'-I}<\bar\rho$ and $0<\abs k_1\le\bar K$. As $\bar\omega_0(I)\in\mathcal D^1_{\bar\gamma}$, $\abs{k\cdot\bar\omega_0(I)}\ge\bar\gamma\bar K^{-\tau}$; the mean value inequality on the segment $[I,I']\subset(\overline V)_{\rho_c}$ gives $\abs{k\cdot(\tilde\omega_0(I')-\bar\omega_0(I))}\le\abs k_1\bar M_0\bar\rho\le\bar K\bar M_0\bar\rho=\bar\gamma\bar K^{-\tau}/2$. So $\abs{k\cdot\tilde\omega_0(I')}\ge\bar\gamma/(2\bar K^\tau)=\alpha_{\rm av}$ (CP10 (39)).
\proofstep{(R2) and the remainder}\label{stepP:p5:av-R2} $\alpha_{\rm av}d=\bar\gamma^2s_*/(4\bar M_0\bar K^{2\tau+1})\ge G^2s_*\mu_2r^{-5}/(4\bar M_0)$ (first branch of $\bar\gamma$). (R2): $2^{12}(1+b)\bar Ks_*\mu_2\le\alpha_{\rm av}d$ holds if $G^2\ge2^{14}(1+b)\bar M_0$ (as $\bar Kr^5\le1$); and $400(1+b)\mu_2/(\alpha_{\rm av}d)\le r^5$ holds if $G^2\ge1600(1+b)\bar M_0/s_*$. Both are conditions on $G$ (Lemma~\ref{bpP:p5:kam-cond}). Then $\wn{\widetilde P_{\rm rem}}{}{}\le\mu_2r^5+\mu_2r^5=2\Lambda_P\mu r^5$ (\Pref{bpP:p2:averaged-remainder}).
\proofstep{Identification}\label{stepP:p5:av-id} In \Pref{bpP:p2:averaged-remainder}, $\mu_2\langle P_2\rangle=\mu\tilde P_{\rm av}=\mu\tilde N+\mu Q$ (Lemma~\ref{bpP:p5:birkhoff-ext}), which gives the displayed identity. \textbf{Displacements:} \Pref{bpP:p2:use5}\,(3) converts \Pref{bpP:p2:averaged-remainder}'s bounds to $\alpha_{\rm av}s_*\abs{\Delta I}\le6\mu_2$, $\alpha_{\rm av}\varepsilon_{\rm an}\abs{\Delta p},\abs{\Delta q}\le9\mu_2$; insert $\alpha_{\rm av}=\bar\gamma/2\bar K^\tau$ and $\bar\gamma\ge G\sqrt{\mu_2}\bar K^{\tau+1/2}r^{-5/2}$: $\mu_2\bar K^\tau/\bar\gamma\le\sqrt{\mu_2}\,r^{5/2}\bar K^{-1/2}/G\le r^{11/2}/G$.
\end{proof}

\subsection{The KAM Hamiltonian in oscillator action--angle variables}\label{bpP:p5:kam-ham}

\begin{definition}[Polar chart, composite chart $\Phi$]\label{bpP:p5:phi}
$\mathrm{Pol}(\psi,J)=\bigl((\varphi,q),(I,p)\bigr)$ with $\varphi=\psi_{\rm f}$, $I=J_1$, $(q_j,p_j)=\sqrt{2J_{2,j}}(\sin\psi_{{\rm s},j},\cos\psi_{{\rm s},j})$ ($J_2\in(0,\infty)^b$): the map $\Pi$ of \Pref{bpP:p4:s-disc}, in the frozen order (position, momentum); it extends holomorphically to $\{\operatorname{Re}J_2>0\}\times\C^n$ (principal square root). Let $O_0:=\{(\psi,J)\in\Ph n:\ J_1\in\operatorname{Re}(\check D_{\bar\rho/2}),\ J_2\in(\lambda_1/2,2\lambda_2)^b\}$ (open, invariant under $\psi\mapsto\psi+2\pi k$) and $\Phi:=\bar\Psi\circ\mathrm{Pol}$ (real restriction on $O_0$, any extension elsewhere), $\Phi:O_0\to\Ph n$.
\end{definition}

\begin{lemma}[$\Phi$ is a $1$-chart with the lattice properties of \Pref{bpP:p1:cyl-pipeline}]\label{bpP:p5:phi-chart}
$\Phi$ is a symplectic chart on $O_0$ (Definition~\ref{bpP:p1:chart}, $c=1$), $O_0+(2\pi k,0)=O_0$, $\Phi(z+(2\pi k,0))=\Phi(z)+\operatorname{shift}(k_{\rm f})$ for $k\in\Z^n$, $z\in O_0$, and $\Phi(z')=\Phi(z)+\operatorname{shift}(k_{\rm f})\Rightarrow z'=z+(2\pi m,0)$ for some $m\in\Z^n$. Moreover $H_\mu\circ\Phi=h_\mu+f_\mu$ on $O_0$, where, with $N(J)=\tilde N(J_1,J_2)$,
\[h_\mu(J)=\tilde H_0(J_1)+\mu\tilde N(J),\qquad f_\mu=\bigl(\mu Q+\widetilde P_{\rm rem}\bigr)\circ\mathrm{Pol}.\]
\end{lemma}
\begin{proof}
\proofstep{Image}\label{stepP:p5:pc-img} For $z\in O_0$: $\abs{\bar p},\abs{\bar q}\le\sqrt{2J_2}<2\sqrt{\lambda_2}\le2r<\varepsilon_{\rm an}/2$ and $I\in\check D_{\bar\rho/2}$: $\mathrm{Pol}(O_0)\subset\mathcal W\cap\R^{2n}$, so $\bar\Psi$ is defined and holomorphic there. $\bar\Psi$ real analytic and symplectic: \Pref{bpP:p1:hol-chart} (with Proposition~\ref{bpP:p5:avg}); $\mathrm{Pol}$ is a $1$-chart on $\{J_2>0\}$: \Pref{bpP:p1:typical}(1) for each oscillator and the identity on $(\varphi,I)$; composition: \Pref{bpP:p1:typical}(3).
\proofstep{Lattice}\label{stepP:p5:pc-lat} $\mathrm{Pol}(\psi+2\pi k,J)=\mathrm{Pol}(\psi,J)+\operatorname{shift}(k_{\rm f})$ ($\psi_{\rm s}$ enters through $2\pi$-periodic functions; $\operatorname{shift}$ moves only the fast angles) and $\bar\Psi(w+\operatorname{shift}(k_{\rm f}))=\bar\Psi(w)+\operatorname{shift}(k_{\rm f})$ (Proposition~\ref{bpP:p5:avg}: $\bar\Psi-\mathrm{id}$ periodic). Injectivity modulo the lattice: if $\Phi(z')=\Phi(z)+\operatorname{shift}(k_{\rm f})=\Phi(z+(2\pi(k_{\rm f},0),0))$, injectivity of $\bar\Psi$ gives $\mathrm{Pol}(z')=\mathrm{Pol}(z+(2\pi(k_{\rm f},0),0))$, and $\mathrm{Pol}$ is injective modulo $\psi\mapsto\psi+2\pi k$ on $\{J_2>0\}$ (\Pref{bpP:p4:disc-measure}(a)).
\proofstep{Hamiltonian}\label{stepP:p5:pc-ham} $H_\mu\circ\bar\Psi$ is given by Proposition~\ref{bpP:p5:avg}; compose with $\mathrm{Pol}$, $R_j\circ\mathrm{Pol}=J_{2,j}$ (\Pref{bpP:p0:oscaction}).
\end{proof}

\begin{lemma}[Hypotheses of the two-scale KAM theorem]\label{bpP:p5:kam-hyp}
Let $H_3(J,\theta):=\tilde H(J,2\pi\theta)/2\pi$ where $\tilde H=h_\mu+f_\mu$ in $(J,\psi)$, i.e. $h_3=h_\mu/2\pi$, $f_3=f_\mu(\cdot,2\pi\theta)/2\pi$. There is $C_1\ge1$, depending on the data only (Remark~\ref{bpP:p5:order}, items 2--3), such that for $0<r<r_*$, $0<\mu<\min\{\mu_0,r^6/(\log\frac1r)^c\}$, the hypotheses of \Pref{bpP:p3:two-scale-kam} hold with: $n_1=a$, $n_2=b$, $\tau=n+1$, $D=\check D\times A_{\rm sh}$, $\rho=\lambda_1/2$, $A=D_\rho$, $s,\bar s$ of Definition~\ref{bpP:p5:params}, and
\[M=C_1,\quad\hat M=C_1\mu,\quad\bar M=\bar M_2=C_1\mu^{-1},\quad\bar M_1=C_1,\quad E=C_1\mu r^5\qquad(\text{the values of \Pref{bpP:p3:interface}\,(2)}).\]
Explicitly $C_1:=\max\{1,\ 2\pi C_U,\ C_R\coth(\sigma_{\rm av}/4)^n/2\pi\}$ with $C_R:=2^5C_Q+2\Lambda_P$ and $C_U$ of Lemma~\ref{bpP:p5:freq-inj}.
\end{lemma}
\begin{proof}
\proofstep{Domain}\label{stepP:p5:kh-dom} $D_\rho=\check D_\rho\times(A_{\rm sh})_\rho$ (max norm) $\subset\mathcal E_1$, since $\rho\le\bar\rho/2\le\rho_c$ and $\abs{J_2}<\lambda_2'+\rho<1$; also $\operatorname{Re}J_2>\lambda_1/2$. The polar image of $D_\rho\times\{\abs{\operatorname{Im}\psi}<\sigma_{\rm av}\}$: $\abs{(\bar p,\bar q)}_\infty\le\sqrt{2\abs{J_2}}\,e^{\sigma_{\rm av}}<2r\le\varepsilon_{\rm an}/2$ (as $\abs{J_2}<\lambda_2$, $\sigma_{\rm av}\le\frac1{12}$), $\abs{\operatorname{Im}\varphi}<\sigma_{\rm av}$, $J_1\in\check D_{\bar\rho/2}$: inside $\mathcal W$. So $f_\mu$ is holomorphic on $D_\rho\times\{\abs{\operatorname{Im}\psi}<\sigma_{\rm av}\}$, $2\pi$-periodic, real on real points; $h_\mu$ is holomorphic on $D_\rho$.
\proofstep{$h$}\label{stepP:p5:kh-h} Lemma~\ref{bpP:p5:freq-inj}: $\partial h_3=F_\mu/2\pi$ injective, $U_3=U/2\pi$ nonsingular on $A$, $\norm{U_3}\le C_U$, $\norm{\hat U_3}\le C_U\mu$, $\norm{U_3^{-1}}\le2\pi C_U\mu^{-1}$, $\norm{T_1}\le2\pi C_U$, $\norm{T_2}\le2\pi C_U/\mu$ (the factors $2\pi$ from $U_3^{-1}=2\pi U^{-1}$). $\mu<\mu_*$ by Lemma~\ref{bpP:p5:param-facts}(a).
\proofstep{$f$}\label{stepP:p5:kh-f} On the polar image $\abs x_\infty<2r\le r_1$: $\abs{\mu Q}\le\mu C_Q(2r)^5$ (Proposition~\ref{bpP:p5:uniform-remainder}; $I\in\hat D_c'$), $\abs{\widetilde P_{\rm rem}}\le\wn{\widetilde P_{\rm rem}}{}{}\le2\Lambda_P\mu r^5$ (\Nref{bpN:p4:sup-weighted}); so $\sup\abs{f_\mu}\le C_R\mu r^5$ on the strip $\sigma_{\rm av}$. Sup-to-weighted conversion (proof of Definition~\ref{bpP:p5:consts}, all $n$ angles, loss $\sigma_{\rm av}-\sigma_{\rm tot}=\sigma_{\rm av}/2$): $\norm{f_3}_{\rho,\sigma_{\rm tot}}\le C_R\coth(\sigma_{\rm av}/4)^n\mu r^5/2\pi\le C_1\mu r^5=E$.
\proofstep{Normalizations}\label{stepP:p5:kh-norm} $\hat M=C_1\mu\le M=C_1$, $\bar M_i\le\bar M$ since $\mu\le1$. All numerical constants are upper bounds as required (\Pref{bpP:p3:rem-norm}).
\end{proof}

\subsection{The scale constants, the KAM condition, and $r_K$ (replacing \Pref{bpP:p4:use5})}\label{bpP:p5:scales}

\begin{lemma}[Scale constants and the KAM condition]\label{bpP:p5:kam-cond}
Let $\hat c=\hat c(n,\tau)=\max\{2^{13}n,24^{\tau+1}/3\}$ (\Pref{bpP:p3:constants}) and $C_1,s$ as above. There are constants, chosen \emph{in this order}:
\begin{enumerate}
\item $\kappa_0\ge C_1^2$ with $\kappa_0^2>9\hat cC_1^4K_0^{2\tau+2}$, $K_0=\frac6s\log_+(\kappa_0^2/C_1^4)$ (exists: $\kappa_0^2/(\log\kappa_0)^{2\tau+2}\to\infty$);
\item $\kappa_1$ with $\kappa_1^2>4\hat cC_1^2$;
\item $G\ge\max\{2\bar M_0\kappa_1,\ (2^{14}(1+b)\bar M_0)^{1/2},\ (1600(1+b)\bar M_0/s_*)^{1/2}\}$.
\end{enumerate}
With them, for all $r\in(0,1)$ small enough that $\rho=\lambda_1/2$ makes sense and all $\mu\in(0,1]$ (Part~3's formulas, with $E=C_1\mu r^5$ etc. of Lemma~\ref{bpP:p5:kam-hyp}): $L=C_1/\mu$, $x=C_1^4r^5/\gamma_1^2=C_1^4/\kappa_0^2$, $K=K_0$ \emph{independent of $r,\mu$},
\[\hat\rho=\min\Bigl\{\frac{\kappa_0r^{5/2}}{3C_1K_0^{\tau+1}},\ \frac{\kappa_1r^{5/2}}{2}\Bigr\},\qquad\hat E=\frac{EL}{\hat\rho^2}=\frac{C_1^2r^5}{\hat\rho^2}=\max\Bigl\{\frac{9C_1^4K_0^{2\tau+2}}{\kappa_0^2},\frac{4C_1^2}{\kappa_1^2}\Bigr\}<\frac1{\hat c},\]
that is, the \textbf{KAM condition} $\hat c\hat E<1$ holds, $\gamma_1\ge\gamma_2$, and $\tau>n$.
\end{lemma}
\begin{proof}
\proofstep{Evaluation}\label{stepP:p5:kc-eval} $L=\max\{\bar M,M^{-1},\hat M^{-1}\}=\max\{C_1/\mu,1/C_1,1/(C_1\mu)\}=C_1/\mu$ ($C_1\ge1$, $\mu\le1$). $EM^2L=C_1^4r^5$, so $x=C_1^4r^5/\gamma_1^2=C_1^4/\kappa_0^2\le1$ and $K=\frac6s\log_+(x^{-1})=K_0$ exactly (\Pref{bpP:p3:interface}\,(2)). Since $\gamma_2/(3\hat MK^{\tau+1})=\mu\kappa_0r^{5/2}/(3C_1\mu K_0^{\tau+1})=\gamma_1/(3MK^{\tau+1})$, $\hat\rho$ is as displayed; $\hat E=EL/\hat\rho^2=C_1^2r^5/\hat\rho^2$; $\hat\rho^{-2}r^5=\max\{(3C_1K_0^{\tau+1}/\kappa_0)^2,4/\kappa_1^2\}$.
\proofstep{Smallness}\label{stepP:p5:kc-small} Items 1, 2 give $\hat c\hat E<1$. $\gamma_2=\mu\gamma_1\le\gamma_1$.
\end{proof}

\begin{lemma}[The radius $r_K$ (what Part~4 should state instead of $r_K\le C\hat\rho\le C\gamma_2$)]\label{bpP:p5:rK}
In the situation of Lemma~\ref{bpP:p5:kam-cond}, $r_K:=32\hat E\hat\rho$ (the enlargement radius of \Pref{bpP:p3:two-scale-kam}) satisfies, for all admissible $(r,\mu)$,
\[r_K<\tfrac{1}{256n}\hat\rho\le\tfrac1{512n}\lambda_1,\qquad r_K<\min\{t,\lambda_1\},\qquad r_K<\tfrac{\bar\rho}2,\qquad\frac{2r_K}{\ell}\le\frac{\kappa_1}{64n}\sqrt r,\qquad\Bigl(1+\frac{2r_K}\ell\Bigr)^b<2.\]
\end{lemma}
\begin{proof}
\proofstep{Direct}\label{stepP:p5:rk-dir} $\hat c\ge2^{13}n$ and $\hat c\hat E<1$ give $\hat E<2^{-13}/n$, so $r_K=32\hat E\hat\rho<\hat\rho/(256n)$, and $\hat\rho\le\rho=\lambda_1/2$. Further $\lambda_1=\kappa_1r^{5/2}<\sqrt r=t$ for $r<1$ ($\kappa_1r^2<1$ by (r4)), and $\lambda_1\le\bar\rho$ (Lemma~\ref{bpP:p5:param-facts}(b)). Finally $\ell\ge r^2/4$ gives $2r_K/\ell\le2\kappa_1r^{5/2}\cdot4/(512n\,r^2)=\kappa_1\sqrt r/(64n)$, and $(1+x)^b\le e^{bx}$ with $x=2r_K/\ell$, $b\le n$, $\kappa_1\sqrt r\le\frac1{12}$ (r4): $bx\le1/768$, so $(1+x)^b<2$.
\end{proof}
\begin{remark}[Interface with Part~4 (reviews P4-01, P4-02)]\label{bpP:p5:use5-fix}
Part~4's \Pref{bpP:p4:use5} and \Pref{bpP:p4:domain-collar} have been corrected in the meantime (chain $r_K\le C\hat\rho\le C\min\{\gamma_1,\gamma_2/\mu\}$; unnormalised collar bound with a positive-volume premise). Part~5 does not use the $\bar D_t$-version of \Pref{bpP:p4:domain-collar} as a black box: for the Cantor-type domain $\check D\times A_{\rm sh}$ it re-derives the corresponding bound in Lemma~\ref{bpP:p5:budget}, Step~\ref{stepP:p5:bu-collar} (the same Steps~\ref{stepP:p4:dc-prod}--\ref{stepP:p4:dc-diff}, with the extra piece $\bar D_t\setminus\check D$), and gets $r_K$ from Lemma~\ref{bpP:p5:rK}. Positive volume is ensured by $A\sqrt r\le\frac12\meas V$ (Lemma~\ref{bpP:p5:param-facts}(c)).
\end{remark}

\subsection{Location, the tori, and the Kolmogorov set}\label{bpP:p5:tori}

\begin{lemma}[Location of the tori in the cylinder]\label{bpP:p5:location}
Let $\varphi_\omega=(v,\vartheta+u)$ be a torus of \Pref{bpP:p3:two-scale-kam} (so $v\in\operatorname{Re}(D_{r_K})$), and let $z=\Phi(2\pi(\vartheta+u),v)\in\Ph n$. Then $v\in O_0$-range ($v_1\in\operatorname{Re}(\check D_{\bar\rho/2})$, $v_2\in(\lambda_1/2,2\lambda_2)^b$), $I(z)\in V$ and $p_j^2+q_j^2<r^2$ for all $j$; thus $z\in\mathcal C(V,r)$. Here $\Delta_{\max}$ is that of Proposition~\ref{bpP:p5:avg}, and (r5) is the condition $\sqrt2\Delta_{\max}\le\kappa_1r^{3/2}$ (satisfied for small $r$ since $\Delta_{\max}\le\frac{18}{G\varepsilon_{\rm an}}r^{11/2}$).
\end{lemma}
\begin{proof}
\proofstep{Range}\label{stepP:p5:lo-range} $v\in\operatorname{Re}(D_{r_K})$ means $\operatorname{dist}(v_1,\check D)<r_K<\bar\rho/2$ (Lemma~\ref{bpP:p5:rK}) and $v_2\in(\lambda_1-r_K,\lambda_2'+r_K)^b\subset(\lambda_1/2,2\lambda_2)^b$.
\proofstep{Actions}\label{stepP:p5:lo-act} $I=\bar I+(I-\bar I)$ with $\bar I=v_1$, $\operatorname{dist}(\bar I,\bar D_t)<r_K$ and $\abs{I-\bar I}\le\Delta_I$ (Proposition~\ref{bpP:p5:avg}); $r_K+\Delta_I<t$, hence $I\in N_{r_K+\Delta_I}(\bar D_t)\subset V$ by \Pref{bpP:p4:collar}(c).
\proofstep{Discs}\label{stepP:p5:lo-disc} For the polar point $(\bar p,\bar q)$: $\bar p_j^2+\bar q_j^2=2v_{2,j}<2(\lambda_2'+r_K)=r^2-4\lambda_1+2r_K\le r^2-3\lambda_1$ (Lemma~\ref{bpP:p5:rK}), so $\sqrt{\bar p_j^2+\bar q_j^2}<r-\frac{3\lambda_1}{2r}=r-\frac32\kappa_1r^{3/2}$. With $\abs{p-\bar p},\abs{q-\bar q}\le\Delta_{\max}$: $\sqrt{p_j^2+q_j^2}\le\sqrt{\bar p_j^2+\bar q_j^2}+\sqrt2\Delta_{\max}<r-\frac32\kappa_1r^{3/2}+\kappa_1r^{3/2}<r$.
\end{proof}

\begin{definition}[Torus family and Kolmogorov set]\label{bpP:p5:kset}
Let $\Omega_*:=\omega_0(D)\cap\mathcal D_{\gamma_1,\gamma_2,\tau}$ ($\omega_0=\partial h_3=F_\mu/2\pi$), $D_0:=D^{\rm dio}=\omega_0^{-1}(\Omega_*)\cap D$, and for each Borel $B\subset D_0$ with $\Omega_B:=\omega_0(B)$ and $\varphi_\omega$ as in \Pref{bpP:p3:two-scale-kam},
\[K_\omega(\theta):=\Phi\bigl(S_{2\pi}\,\pi\,\varphi_\omega(\theta)\bigr)\quad(\omega\in\Omega_*),\qquad S_B:=\bigcup_{\omega\in\Omega_B}K_\omega(\R^n),\qquad\mathsf K_B:=\bigcup_{\omega\in\Omega_B}\varphi_\omega(\R^n)\subset\R^n\times\R^n\ \text{(coordinates }(J,\vartheta)).\]
($\pi(J,\vartheta)=(\vartheta,J)$, $S_{2\pi}$ of \Pref{bpP:p1:period-convert}.) $D_0$ is relatively closed in $D$ ($\mathcal D_{\gamma_1,\gamma_2,\tau}$ closed, $F_\mu$ continuous), $D=\check D\times A_{\rm sh}$ is $\sigma$-compact, so $D_0$ is $\sigma$-compact.
\end{definition}

\begin{proposition}[Each torus is a frozen cylinder torus with frozen frequencies]\label{bpP:p5:torus-conv}
For $0<r<r_*$ and admissible $\mu$, $\omega\in\Omega_*$, $w:=\omega$: $\operatorname{IsCylTorus}(H_\mu,\mathcal C(V,r),w,K_\omega)$; $w$ is $(\min\{\gamma_1,\gamma_2\},\tau)$-Diophantine with $\tau=n+1>a+b$; $\abs{w^{\rm f}}_\infty\le C_w$, $\abs{w^{\rm s}}_\infty\le C_w\mu$ with $C_w:=\frac{M_\omega+A_N}{2\pi}+1$, $M_\omega:=\sup_{\overline V}\abs{\bar\omega_0}$; $K_\omega(\R^n)\subset\mathcal C(V,r)$.
\end{proposition}
\begin{proof}
\proofstep{Chart data}\label{stepP:p5:tc-chart} Lemma~\ref{bpP:p5:phi-chart} supplies $O_0$, $\Phi$ with the hypotheses of \Pref{bpP:p1:cyl-pipeline} ($d=n$); $G_3=(2\pi)^{-1}(H_\mu\circ\Phi)(2\pi z_1,z_2)$, which by Lemma~\ref{bpP:p5:phi-chart} equals $H_3$ in the frozen order $\pi$ on $O=S_{2\pi}^{-1}(O_0)$.
\proofstep{$\operatorname{LT}(G_3,O,w,K_3)$, $K_3=\pi\circ\varphi_\omega$}\label{stepP:p5:tc-LT} (L1): $\varphi_\omega$ is real analytic (\Pref{bpP:p3:two-scale-kam}(a)); $K_3(\R^n)\subset O$ by Lemma~\ref{bpP:p5:location}, Step~\ref{stepP:p5:lo-range}; $G_3$ is differentiable at $K_3(\theta)$ (chart-field Lemma~\ref{bpP:p1:chart-field}: $H_\mu$ is analytic near $\Phi(\ldots)\in\mathcal C(V,r)$ since $r<r_0$, (H1), $H_0$ analytic on $W$). (L2): $v,u$ are $1$-periodic. (L3): injectivity modulo translations (Proposition~\ref{bpP:p3:two-scale-kam}, with equal $\omega$). (L4), (L5): Proposition~\ref{bpP:p3:two-scale-kam}(c); the displayed vanishing sum is $\Omega(DK_3\,\cdot,DK_3\,\cdot)$ up to sign. (L6): Proposition~\ref{bpP:p3:two-scale-kam}(b), with $\pi$ (\Pref{bpP:p1:order}), where $H_3=G_3$ on the open set $\operatorname{Re}(D_\rho)\times\R^n\ni K_3(\theta)$ ($r_K<\rho$) so the fields agree.
\proofstep{Pipeline}\label{stepP:p5:tc-pipe} \Pref{bpP:p1:cyl-pipeline} with $U=\mathcal C(V,r)\supset\Psi(K_3(\R^n))=K_\omega(\R^n)$ (Lemma~\ref{bpP:p5:location}) gives $\operatorname{IsCylTorus}(H_\mu,\mathcal C(V,r),w,K_\omega)$ with the same $w=\omega$.
\proofstep{Frequencies}\label{stepP:p5:tc-freq} $\omega=\omega_0(I_0)$, $I_0=(J_1,J_2)\in D$ real, $\abs{J_2}\le\lambda_2\le1$: $w^{\rm f}=(\bar\omega_0(J_1)+\mu G)/2\pi$, $w^{\rm s}=\mu\Theta/2\pi$, bounded using $\abs{G},\abs\Theta\le A_N$. Diophantine: \Pref{bpP:p1:dio-sources}(a) ($\gamma_1,\gamma_2>0$, $\tau>n$).
\end{proof}

\begin{lemma}[Borelness, whole tori, cell measure]\label{bpP:p5:cell-transfer}
Let $B\subset D_0$ be $\sigma$-compact and $\mathsf A_B:=\mathsf K_B\cap(\R^n\times[0,1)^n)$ (coordinates $(J,\vartheta)$; the \emph{measured representative}, an \emph{output}-angle cell, not the image of the parameter cell). (a) $\mathsf A_B$ is $\sigma$-compact; $S_B$ is Borel, $S_B\subset\mathcal C(V,r)$, and every point of $S_B$ lies on some $K_\omega(\R^n)\subset S_B$. (b) $\meas(S_B\cap\mathcal F(V,r))=(2\pi)^n\meas\mathsf A_B$.
\end{lemma}
\begin{proof}
\proofstep{$\sigma$-compactness}\label{stepP:p5:ct-cell} Let $\mathsf K^{\rm c}_B:=\Phi_3(\check\imath(B)\times[0,1)^n)$ (notation of Part~3), the image of the \emph{parameter} cell. $\check\imath$ and $\Phi_3$ are continuous (\Pref{bpP:p3:limit-actions}, \Pref{bpP:p3:limit-phi}) and $B\times[0,1)^n$ is $\sigma$-compact, so $\mathsf K^{\rm c}_B$ is $\sigma$-compact. Since $\varphi_\omega(\vartheta+k)=\varphi_\omega(\vartheta)+(0,k)$, $\mathsf K_B=\bigcup_{k\in\Z^n}(\mathsf K^{\rm c}_B+(0,k))$ is $\sigma$-compact and invariant under $(0,k)$. $\mathsf A_B$ is the intersection of a $\sigma$-compact set with an $F_\sigma$ set, hence $\sigma$-compact. (The output angles of $\mathsf K^{\rm c}_B$ are $\vartheta+u$ and need not lie in $[0,1)^n$, so $\mathsf K^{\rm c}_B\ne\mathsf A_B$ in general; only $\mathsf A_B$ is measured below.)
\proofstep{Decomposition}\label{stepP:p5:ct-dec} Put $T:=\Phi(S_{2\pi}\pi\,\mathsf A_B)$; $S_{2\pi}\pi\mathsf A_B\subset Q'':=\{\psi\in[0,2\pi)^n\}\cap O_0$ (Lemma~\ref{bpP:p5:location}), $\sigma$-compact, and $\Phi$ is continuous on $O_0$, so $T$ is $\sigma$-compact. If $(J,\vartheta)=\varphi_\omega(\theta)\in\mathsf K_B$ write $\vartheta=\vartheta'+m$, $\vartheta'\in[0,1)^n$, $m\in\Z^n$: $(J,\vartheta')\in\mathsf A_B$ and, by Lemma~\ref{bpP:p5:phi-chart}, $K_\omega(\theta)=\Phi(2\pi\vartheta'+2\pi m,J)=\Phi(2\pi\vartheta',J)+\operatorname{shift}(m_{\rm f})\in T+\operatorname{shift}(m_{\rm f})$. Conversely every point of $T$ is some $K_\omega(\theta)$, and $S_B$ is invariant under $\operatorname{shift}(k)$ ($K_\omega(\theta+k)=K_\omega(\theta)+\operatorname{shift}(k_{\rm f})$). So $S_B=\bigcup_{k\in\Z^a}(T+\operatorname{shift}(k))$, Borel; (a) follows with Lemma~\ref{bpP:p5:location}.
\proofstep{Disjointness and injectivity}\label{stepP:p5:ct-disj} If $z,z'\in Q''$ and $\Phi(z)+\operatorname{shift}(k)=\Phi(z')+\operatorname{shift}(k')$, then $\Phi(z')=\Phi(z)+\operatorname{shift}(k-k')$ and (Q2, Lemma~\ref{bpP:p5:phi-chart}) $z'=z+(2\pi m,0)$ with $m\in\Z^n$; both lie in $[0,2\pi)^n$, so $m=0$, $z'=z$ and $k=k'$. Thus $\Phi$ is injective on $S_{2\pi}\pi\mathsf A_B$ and the translates $T+\operatorname{shift}(k)$, $k\in\Z^a$, are pairwise disjoint.
\proofstep{Counting}\label{stepP:p5:ct-count} $S_B\subset\mathcal C(V,r)$, so $S_B\cap\mathcal F=S_B\cap\{\varphi\in[0,2\pi)^a\}$ and by translation invariance of $\meas$, $\meas(S_B\cap\{\varphi\in[0,2\pi)^a\})=\sum_k\meas\bigl((T+\operatorname{shift}(k))\cap\{\varphi\in[0,2\pi)^a\}\bigr)=\sum_k\meas\bigl(T\cap\{\varphi\in[-2\pi k,2\pi(1-k))\}\bigr)=\meas T$ (the slabs partition $\R^a$; the averaging map moves the fast angles, but no physical cell is required to be preserved).
\proofstep{Liouville}\label{stepP:p5:ct-liou} $\Phi$ is a $1$-chart, injective on the measurable set $S_{2\pi}\pi\mathsf A_B$ (Step~\ref{stepP:p5:ct-disj}): \Pref{bpP:p1:liouville} gives $\meas T=\meas(S_{2\pi}\pi\mathsf A_B)=(2\pi)^n\meas\mathsf A_B$ ($S_{2\pi}$ multiplies volume by $(2\pi)^n$, \Pref{bpP:p1:period-convert}(a); $\pi$ permutes coordinates).
\end{proof}


\subsection{The measure budget (replacing \Pref{bpP:p4:budget} and \Pref{bpP:p4:cell-budget} for $D=\check D\times A_{\rm sh}$)}\label{bpP:p5:measure}

\begin{lemma}[Budget on the action side]\label{bpP:p5:budget}
There is $C_{\rm b}<\infty$ (data and the constants $\kappa_0,\kappa_1,G$ only) such that for admissible $(r,\mu)$, with $\mathsf K=\mathsf K_{D_0}$ and $c_n=1+(1+2^n)(1+(4/3)^n)$,
\[\meas\bigl((D\times\T^n)\setminus\mathsf K\bigr)\le c_nC_{\rm b}\sqrt r\ \meas D,\qquad\meas D=\meas\check D\,\ell^b>0.\]
\end{lemma}
\begin{proof}
\proofstep{Reading (3.1)}\label{stepP:p5:bu-read} $\mathsf K=K_\Omega$ of \Pref{bpP:p3:two-scale-kam} for $(\Omega_*,r_K)$; $D\subset\operatorname{Re}(D_{r_K})$, so $(D\times\T^n)\setminus\mathsf K\subset(\operatorname{Re}(D_{r_K})\times\T^n)\setminus\mathsf K$, of measure $\le c_n\bigl(\meas(D\setminus D_0)+\meas(\operatorname{Re}(D_{r_K})\setminus D)\bigr)$.
\proofstep{Diophantine part}\label{stepP:p5:bu-dio} $D\setminus D_0\subset\mathcal R:=\{(I,\xi)\in\bar D_t\times A_{\rm sh}:F_\mu(I,\xi)\notin\mathcal D_{\alpha_1,\alpha_2,\tau}\}$, $\alpha_i=2\pi\gamma_i$ (dictionary \Pref{bpP:p4:freq}). The standing data (S4) of \Pref{bpP:p4:data} hold with $\bar D=\bar D_t$ (compact, $\meas>0$ by Lemma~\ref{bpP:p5:param-facts}(c)), $\lambda_0=1$, $\mu_1$ of (r1), and the constants of Lemma~\ref{bpP:p5:freq-data}. \Pref{bpP:p4:two-scale-diophantine-measure} with $\lambda=\lambda_2'\in[\lambda_2/2,1]$, $Q=A_{\rm sh}$, $\bar\gamma_{\rm P4}=\alpha_2/\mu=2\pi\kappa_0r^{5/2}$: $\meas\mathcal R\le[c_1\alpha_1(1+\mu^{\tau-n}/\lambda)+c_2\bar\gamma_{\rm P4}/\lambda]\meas\bar D_t(2\lambda)^b$. Here $c_1,c_2$ (which depend on $\bar D$ only through $1/\meas\bar D$) are $\le\bar c_1,\bar c_2$, data-dependent, as $\meas\bar D_t\ge\frac12\meas V$; $\tau-n=1$, $\mu\le r^6$, $\lambda\ge r^2/4$ give $\mu/\lambda\le4$ and $\bar\gamma_{\rm P4}/\lambda\le8\pi\kappa_0r^{1/2}$. Also $(2\lambda)^b\le(4\ell)^b$ ($\lambda_1\le\lambda_2'/2$), $\meas\bar D_t\le4\meas\check D$. So $\meas(D\setminus D_0)\le4^{b+1}2\pi\kappa_0(5\bar c_1+4\bar c_2)\sqrt r\,\meas D$.
\proofstep{Collar part}\label{stepP:p5:bu-collar} $\operatorname{Re}(D_{r_K})=N_{r_K}(\check D)\times(\lambda_1-r_K,\lambda_2'+r_K)^b$ (max norm, \Pref{bpP:p4:domain-collar}, Step~\ref{stepP:p4:dc-prod}). $N_{r_K}(\check D)\setminus\check D\subset(N_{r_K}(\bar D_t)\setminus\bar D_t)\cup(\bar D_t\setminus\check D)$, of measure $\le At+c_{\rm dio}\bar\gamma\le(A+c_{\rm dio}C_\gamma)\sqrt r$ (\Pref{bpP:p4:collar}(c) with $r_K<t<d$; Lemma~\ref{bpP:p5:dio-one}). As in \Pref{bpP:p4:domain-collar}, Step~\ref{stepP:p4:dc-diff}, and with $\meas\check D\ge\frac14\meas V$ and $(1+2r_K/\ell)^b<2$ (Lemma~\ref{bpP:p5:rK}): $\meas(\operatorname{Re}(D_{r_K})\setminus D)\le\meas D\,[8(A+c_{\rm dio}C_\gamma)/\meas V+b\kappa_1/16]\sqrt r$. (The denominator $\meas\check D>0$ is part of the statement: review P4-01.)
\end{proof}

\begin{proposition}[Lower bound for $\meas(S\cap\mathcal F)$]\label{bpP:p5:lower}
Let $C_*:=\dfrac{A+c_{\rm dio}C_\gamma}{\meas V}+6b\kappa_1+c_nC_{\rm b}$. For $0<r<r_*$ and admissible $\mu$ the set $S:=S_{D_0}$ satisfies $\meas(\mathcal F(V,r)\setminus S)\le C_*\sqrt r\,\meas\mathcal F(V,r)$, i.e. $\meas(S\cap\mathcal F)\ge(1-C_*\sqrt r)\meas\mathcal F$.
\source{CP10 \S3, (55)--(59) and \S4, p.~564: $\meas(P\setminus K_*)\le C(\bar\gamma+\gamma_1+\bar\gamma_2/\epsilon^2+\epsilon^{n_2/2})\meas P$, here with the shell loss $6b\kappa_1\sqrt r$ (Part~4 erratum E4) and the collar term from \texttt{hboundary}.}
\end{proposition}
\begin{proof}
\proofstep{Set up}\label{stepP:p5:lo-setup} By Lemma~\ref{bpP:p5:cell-transfer}(b), $\meas(S\cap\mathcal F)=(2\pi)^n\meas\mathsf L$, $\mathsf L:=\mathsf K\cap(\R^n\times[0,1)^n)$; by Lemma~\ref{bpP:p5:location} (the $(v_1,v_2)$ part) $\mathsf L\subset V\times(0,\lambda_2)^b\times[0,1)^n=:\mathsf F'$, and $\meas\mathcal F=(2\pi)^n\meas\mathsf F'=(2\pi)^n\meas V\lambda_2^b$ (\Pref{bpP:p4:disc-measure}(b)). So $\meas(\mathcal F\setminus S)=(2\pi)^n\meas(\mathsf F'\setminus\mathsf L)$ (no passage through $\mathsf E$ is needed: the two measures are compared in the $(J,\vartheta)$ coordinates).
\proofstep{Bad actions}\label{stepP:p5:lo-bad} $D\times[0,1)^n\subset\mathsf F'$ and $\meas V\lambda_2^b-\meas\check D\ell^b\le\lambda_2^b(\meas V-\meas\check D)+\meas V\,b\lambda_2^{b-1}\cdot3\lambda_1$; $\meas V-\meas\check D\le At+c_{\rm dio}\bar\gamma$ (\Pref{bpP:p4:collar}(b), Lemma~\ref{bpP:p5:dio-one}); $3\lambda_1/\lambda_2=6\kappa_1\sqrt r$. Divide by $\meas V\lambda_2^b$: $\le[(A+c_{\rm dio}C_\gamma)/\meas V+6b\kappa_1]\sqrt r$.
\proofstep{Lost tori}\label{stepP:p5:lo-lost} $\meas((D\times[0,1)^n)\setminus\mathsf L)=\meas((D\times\T^n)\setminus\mathsf K)\le c_nC_{\rm b}\sqrt r\meas D\le c_nC_{\rm b}\sqrt r\meas\mathsf F'$ ($\mathsf K$ is $\Z^n$-invariant in $\vartheta$, Lemma~\ref{bpP:p5:budget}). Add.
\end{proof}

\subsection{Strictness and positivity (ours)}\label{bpP:p5:strict}

\begin{proposition}[Deleting a small positive-measure family of tori]\label{bpP:p5:strictness}
For $0<r<r_*$ and admissible $\mu$ there is a $\sigma$-compact $B\subset D_0$ such that $S':=S_{D_0\setminus B}$ satisfies
\[0<\meas(S'\cap\mathcal F),\qquad(1-(C_*+\tfrac12)\sqrt r)\meas\mathcal F\le\meas(S'\cap\mathcal F)<\meas\mathcal F.\]
\source{Ours; not in CP10 (12) or CP11 Thm.~11.1, where $\meas P>\meas K$ is asserted without proof.}
\sourcegap{The strict inequality is obtained here from the slack $\tfrac12\sqrt r\meas\mathcal F>0$ in the lower bound and the bi-Lipschitz parametrisation of Part~3.}
\end{proposition}
\begin{proof}
\proofstep{Additivity}\label{stepP:p5:st-add} For $\sigma$-compact $B\subset D_0$ put $m(B):=\meas(S_B\cap\mathcal F)=(2\pi)^n\meas\mathsf A_B$ (Lemma~\ref{bpP:p5:cell-transfer}). The tori $\varphi_\omega(\R^n)$, $\omega\in\Omega_*$, are pairwise disjoint (Proposition~\ref{bpP:p3:two-scale-kam}, injectivity modulo translations), so $m$ is countably additive on $\sigma$-compact subsets of $D_0$.
\proofstep{Slices are null}\label{stepP:p5:st-null} Let $x_1(I)=I_1$ (first coordinate) and $H_s=D_0\cap\{x_1=s\}$ (a null set). $\mathsf K^{\rm c}_{H_s}=\Phi_3(\check\imath(H_s)\times[0,1)^n)$ is the image of $\check\imath(H_s)\times[0,1)^n$ under a Lipschitz map ($\check\imath$ is Lipschitz, $\Phi_3$ Lipschitz in the rescaled norm: \Pref{bpP:p3:limit-actions}, \Pref{bpP:p3:limit-phi}); \Pref{bpP:p1:torus-union} gives $\meas\mathsf K^{\rm c}_{H_s}\le\Lambda^{2n}\meas\check\imath(H_s)=0$ (Lipschitz image of a null set); hence $\mathsf K_{H_s}=\bigcup_k(\mathsf K^{\rm c}_{H_s}+(0,k))$ and $\mathsf A_{H_s}\subset\mathsf K_{H_s}$ are null, i.e.\ $m(H_s)=0$.
\proofstep{Intermediate value}\label{stepP:p5:st-ivt} $f(s):=m(D_0\cap\{x_1\le s\})$ is nondecreasing, $f(-\infty)=0$, $f(+\infty)=m(D_0)=\meas(S\cap\mathcal F)\ge(1-C_*\sqrt r)\meas\mathcal F>0$ (Proposition~\ref{bpP:p5:lower}, (r5)). It is right-continuous (decreasing intersections, finite measure, additivity) and left-continuous by Step~\ref{stepP:p5:st-null}. Put $\theta_0:=\min\{\frac12m(D_0),\frac12\sqrt r\meas\mathcal F\}>0$ and choose $s$ with $f(s)=\theta_0$; $B:=D_0\cap\{x_1\le s\}$ is closed in $D_0$, hence $\sigma$-compact, and $D_0\setminus B=D_0\cap\{x_1>s\}$ is $\sigma$-compact.
\proofstep{Conclusion}\label{stepP:p5:st-conc} $\meas(S'\cap\mathcal F)=m(D_0)-\theta_0$. Upper: $\le\meas(S\cap\mathcal F)-\theta_0\le\meas\mathcal F-\theta_0<\meas\mathcal F$ ($S\cap\mathcal F\subset\mathcal F$). Lower: $\ge(1-C_*\sqrt r)\meas\mathcal F-\frac12\sqrt r\meas\mathcal F$; positive since $(C_*+\frac12)\sqrt r\le\frac12$ (r5) and $\meas\mathcal F>0$.
\end{proof}

\subsection{Proof of the frozen theorem}\label{bpP:p5:proof}

\begin{theorem}[P-G; \Pref{bpP:p0:pg}]\label{bpP:p5:main}
\frozen{KAM.Planetary.properlyDegenerateKAM}
Under the hypotheses and data of \Pref{bpP:p0:pg}, with $c:=2n+4$, $C:=\max\{C_*+1,C_w\}$, and $r_*$ of Definition~\ref{bpP:p5:rstar}, the conclusion holds with exactly the printed quantifier order.
\end{theorem}
\begin{proof}
\proofstep{Constants first}\label{stepP:p5:mn-const} Following Remark~\ref{bpP:p5:order}: the data constants of Lemmas~\ref{bpP:p5:ext}--\ref{bpP:p5:freq-inj} (Proposition~\ref{bpP:p5:uniform-remainder}); $C_U,C_1$; $\kappa_0,\kappa_1,G$ (Lemma~\ref{bpP:p5:kam-cond}); $C_\gamma,C_{\rm b},C_*$ (Lemma~\ref{bpP:p5:budget}, Proposition~\ref{bpP:p5:lower}); $C_w$; $r_*$ (Definition~\ref{bpP:p5:rstar}, which is a finite list of conditions on these constants; $r_*<r_0$, $r_*<1$ included); $c=2n+4$ and $C$. None depends on $r$ or $\mu$, and $0<r_*<r_0$, $r_*<1$, $C>0$, $c>0$.
\proofstep{Given $(r,\mu)$}\label{stepP:p5:mn-given} Let $0<r<r_*$, $0<\mu<\mu_0$, $\mu<r^6/(\log\frac1r)^c$ (positive base: $r<1$). All objects of Definition~\ref{bpP:p5:params} are defined, Lemma~\ref{bpP:p5:param-facts} holds; Propositions~\ref{bpP:p5:avg}, Lemmas~\ref{bpP:p5:kam-hyp}, \ref{bpP:p5:kam-cond} verify the hypotheses of \Pref{bpP:p3:two-scale-kam} (including the complex injectivity, Lemma~\ref{bpP:p5:freq-inj}).
\proofstep{The set}\label{stepP:p5:mn-set} $S:=S'$ of Proposition~\ref{bpP:p5:strictness}. By Lemma~\ref{bpP:p5:cell-transfer}(a) it is Borel (so Lebesgue measurable), $S\subset\mathcal C(V,r)$, and every $z\in S$ is $K_\omega(\theta)$ for some $\omega\in\Omega_*$, $\theta\in\R^n$, with $K_\omega(\R^n)\subset S$; Proposition~\ref{bpP:p5:torus-conv} supplies $\operatorname{IsCylTorus}(H_\mu,\mathcal C(V,r),w,K_\omega)$, $\gamma>0$, $\tau>a+b$, $\abs{w^{\rm f}}_\infty\le C_w\le C$, $\abs{w^{\rm s}}_\infty\le C_w\mu\le C\mu$. Hence $\operatorname{IsCylKol}(H_\mu,V,r,\mu,C,S)$ (\Pref{bpP:p0:cylkol}).
\proofstep{Measures}\label{stepP:p5:mn-meas} Proposition~\ref{bpP:p5:strictness}: $0<\meas(S\cap\mathcal F)<\meas\mathcal F$ and $\meas(S\cap\mathcal F)\ge(1-(C_*+\frac12)\sqrt r)\meas\mathcal F>(1-C\sqrt r)\meas\mathcal F$ since $C\ge C_*+1$ and $\meas\mathcal F>0$. If $1-C\sqrt r\le0$, $\operatorname{ofReal}(1-C\sqrt r)=0<\meas(S\cap\mathcal F)$; otherwise $\operatorname{ofReal}(1-C\sqrt r)\meas\mathcal F=(1-C\sqrt r)\meas\mathcal F<\meas(S\cap\mathcal F)$ (\Pref{bpP:p0:totality}). $\meas\mathcal F<\infty$ because $V$ is bounded.
\end{proof}

\begin{remark}[Quantifier audit]\label{bpP:p5:quant}
Frozen order: $\exists\,r_*,C,c\ \forall r,\mu\ \exists S$. Constants chosen before $r$: $r_*,C,c$ and every constant of Remark~\ref{bpP:p5:order}, items 1--5 (all functions of the data). Chosen after $(r,\mu)$: $\bar K,t,\lambda_i,\bar\gamma,\bar\rho,\rho,\gamma_i,\hat\rho,\hat E,r_K$, the set $D$, the family $\varphi_\omega$, the deleted parameters $B$, hence $S$. $K=K_0$ does not depend on $(r,\mu)$ at all. The Diophantine exponent $\tau=n+1$ is data-independent. The frequency bound $C$ and the measure constant coincide, as in the frozen statement.
\end{remark}

\subsection{Errata, deviations, external facts, forward references}\label{bpP:p5:notes}

\paragraph{Suspected defects in the sources.} (1) CP10 p.~563, (60): ``$\sup\abs{\tilde P_{\rm av}}\le C\epsilon^5$'' is asserted from (A$_2''$) (a little-$o$ statement with $\delta$ depending on $(\mu,I,\epsilon)$); supplied as Proposition~\ref{bpP:p5:uniform-remainder} via joint analyticity, periodicity, the jet identification and Schwarz. (2) CP10 p.~563, Step~2: ``we proceed as in Steps 5 and 6'' and the measure estimate (56)--(57) uses $\meas(\operatorname{Re}(V_r)\setminus V)\le C\bar\gamma_2\meas V$, which cannot hold for arbitrary $V$; replaced by \texttt{hboundary}. (3) The neighbourhood of the Cantor set $\check D$ of $\bar\omega_0$-Diophantine actions is not controlled by the collar of $V$; CP10's (57) writes $\operatorname{Re}(D_r)\setminus D\subset(\operatorname{Re}(V_r)\setminus V)\cup(V\setminus D)$ without proof; Lemma~\ref{bpP:p5:budget}, Step~\ref{stepP:p5:bu-collar} proves it. (4) CP10 (20) $\bar\gamma_2=\gamma_*^2\epsilon^{5/2}$ vs (17) lower bound with $(\log(\epsilon^5/\gamma_1^2)^{-1})^{\tau_*+1}$: consistent only because $\gamma_1=\gamma_*^2\epsilon^{5/2}$ makes this logarithm a constant (our $K_0$); not an erratum, but the dependence must be tracked, as done in Lemma~\ref{bpP:p5:kam-cond}. (5) CP10 p.~564 ``replacing $\epsilon$ by $\sqrt{2\check c_2}\epsilon$'': in our frozen normalisation the disc is $\xi_j<r^2/2$ and the shell loss constant is $6b\kappa_1$.
\paragraph{Deviations.} Steps 2--4 of CP10 \S3 are skipped (the hypothesis is already Birkhoff form); the polynomial $\tilde N$ is defined from $\tilde P_{\rm av}$ through its fourth jet; $\check D$ replaces $\bar D\times A$ with an explicit Cantor-set measure estimate; the top of the shell is cut at $\lambda_2'=\lambda_2-2\lambda_1$ so that the tori stay in the open cylinder; $\mu_2=\Lambda_P\mu$ absorbs the sup-norm of $P$ (CP10 normalises $\wn P{}{}=1$).
\paragraph{External facts used (beyond Parts 1--4).} Identity theorem and Cauchy estimates for holomorphic functions of several variables; compactness of closed bounded sets; Cramer's rule; block inversion (Schur complement) and Neumann series; Schwarz lemma in one variable; continuity of measure along monotone sequences; Lebesgue number lemma. Tracks K, N as cited in Parts 1--4.
\paragraph{Forward and cross-track references.} Cited from Part~1: \texttt{cyl-pipeline}, \texttt{hol-chart}, \texttt{chart-field}, \texttt{typical}, \texttt{liouville}, \texttt{torus-union}, \texttt{dio-sources}, \texttt{period-convert}; Part~2: \texttt{averaged-remainder}, \texttt{use5}; Part~3: \texttt{two-scale-kam}, \texttt{interface}, \texttt{period-dictionary}, \texttt{constants}, \texttt{rem-norm}, \texttt{limit-actions}, \texttt{limit-phi}; Part~4: \texttt{collar}, \texttt{lip-inverse}, \texttt{affine-slab}, \texttt{lattice-tail}, \texttt{two-scale-diophantine-measure}, \texttt{data}, \texttt{freq}, \texttt{domain-collar}, \texttt{disc-measure}. Tracks N: \texttt{bpN:p4:norm}, \texttt{bpN:p4:sup-weighted}. Parts~12--13 are not used by P-G. Part~4 must repair \texttt{use5} and \texttt{domain-collar} as stated in Remark~\ref{bpP:p5:use5-fix}.
\paragraph{Label index.} \texttt{bpP:p5:ext}, \texttt{consts}, \texttt{birkhoff-ext}, \texttt{uniform-remainder}, \texttt{lip-complex}, \texttt{bilip}, \texttt{dio-one}, \texttt{freq-data}, \texttt{freq-inj}, \texttt{params}, \texttt{rstar}, \texttt{param-facts}, \texttt{avg}, \texttt{phi}, \texttt{phi-chart}, \texttt{kam-hyp}, \texttt{kam-cond}, \texttt{rK}, \texttt{location}, \texttt{kset}, \texttt{torus-conv}, \texttt{cell-transfer}, \texttt{budget}, \texttt{lower}, \texttt{strictness}, \texttt{main}.
\endgroup
